% ============================================================================
%  SECCIÓN 3.17: CONSTRUCCIÓN DE LAS PANTALLAS
% ============================================================================

\section{Construcci\'on de las pantallas}

Siguiendo la gu\'ia del avance, se construyeron las pantallas principales de la
aplicaci\'on m\'ovil \nombreApp{}, tomando como referencia los wireframes
dise\~nados en el primer avance (secci\'on de wireframes iniciales). A
continuaci\'on se presenta cada pantalla desarrollada, incluyendo su
descripci\'on, los componentes utilizados y una captura de la interfaz
implementada.

\subsection{Pantalla de bienvenida (Splash)}

La pantalla de bienvenida es la primera pantalla que se muestra al usuario al
abrir la aplicaci\'on. Contiene el logotipo de la aplicaci\'on, el nombre
\textit{``Lugares Interactivos''}, un mensaje tagline
\textit{``Descubre, explora y conecta''} y un bot\'on \textit{``Comenzar''}
que redirige al usuario hacia la pantalla de inicio de sesi\'on.

El dise\~no utiliza un fondo con color primario (\texttt{\#2563EB}) y una
forma decorativa curva en la parte superior que le da identidad visual a la
aplicaci\'on. El bot\'on \textit{``Comenzar''} se presenta en color blanco
con sombra para generar contraste y facilitar la interacci\'on.

\begin{figure}[H]
\centering
  \figuraAPA{fig:pantalla-splash}{Pantalla de bienvenida (Splash) de la aplicaci\'on \nombreApp{}}{%
    \includegraphics[width=0.35\linewidth]{pantalla-splash.png}%
  }{Elaboraci\'on propia (captura de pantalla de la aplicaci\'on).}
\end{figure}

Los elementos visuales de esta pantalla se resumen en la tabla.

\begin{table}[H]
  \centering
  \caption{Elementos de la pantalla de bienvenida}
  \contenidoConFuente{%
    \begin{tablaAPA}
      \begin{tabular}{@{}p{3.2cm}p{12.6cm}@{}}
        \toprule
        \textbf{Elemento} & \textbf{Descripci\'on} \\
        \midrule
        Logotipo & Imagen del logo de la aplicaci\'on dentro de un
        contenedor circular blanco con sombra. \\
        \addlinespace
        Nombre de la app & Texto \textit{``Lugares Interactivos''} en tama\~no
        grande, color blanco y peso bold. \\
        \addlinespace
        Tagline & Mensaje breve \textit{``Descubre, explora y conecta''} en
        tama\~no medio y color secundario claro. \\
        \addlinespace
        Bot\'on Comenzar & Bot\'on rectangular de ancho completo con color
        blanco, texto en color primario y sombra suave. \\
        \bottomrule
      \end{tabular}
    \end{tablaAPA}%
  }{Elaboraci\'on propia en base al c\'odigo fuente de la aplicaci\'on.}
\end{table}

\subsection{Pantalla de inicio de sesi\'on}

La pantalla de inicio de sesi\'on permite al usuario acceder a su cuenta
mediante la ingresa de su correo electr\'onico y contrase\~na. Incluye
validaci\'on en tiempo real de los campos, mensajes de error y un enlace hacia
la pantalla de registro para usuarios nuevos.

La interfaz utiliza el componente \texttt{KeyboardAvoidingView} para ajustarse
cuando el teclado virtual aparece, y un \texttt{ScrollView} que permite
desplazarse por el formulario cuando el contenido excede la pantalla.

\begin{figure}[H]
\centering
  \figuraAPA{fig:pantalla-login}{Pantalla de inicio de sesi\'on de la aplicaci\'on \nombreApp{}}{%
    \includegraphics[width=0.35\linewidth]{pantalla-login.png}%
  }{Elaboraci\'on propia (captura de pantalla de la aplicaci\'on).}
\end{figure}

Los elementos de esta pantalla se presentan en la tabla.

\begin{table}[H]
  \centering
  \caption{Elementos de la pantalla de inicio de sesi\'on}
  \contenidoConFuente{%
    \begin{tablaAPA}
      \begin{tabular}{@{}p{3.2cm}p{12.6cm}@{}}
        \toprule
        \textbf{Elemento} & \textbf{Descripci\'on} \\
        \midrule
        Encabezado & T\'itulo \textit{``Bienvenido''} y subt\'itulo
        \textit{``Inicia sesi\'on para continuar''}. \\
        \addlinespace
        Campo email & Campo de texto con validaci\'on de formato de correo
        electr\'onico y mensaje de error si el campo est\'a vac\'io o es
        inv\'alido. \\
        \addlinespace
        Campo contrase\~na & Campo de texto con entrada segura
        (\textit{secureTextEntry}), m\'inimo 6 caracteres. \\
        \addlinespace
        Bot\'on \textit{``Iniciar sesi\'on''} & Bot\'on de color primario que
        muestra un indicador de carga mientras se procesa la solicitud. \\
        \addlinespace
        Enlace \textit{``Registrate''} & Texto enlazado que navega hacia la
        pantalla de registro. \\
        \bottomrule
      \end{tabular}
    \end{tablaAPA}%
  }{Elaboraci\'on propia en base al c\'odigo fuente de la aplicaci\'on.}
\end{table}

\subsection{Pantalla de registro}

La pantalla de registro permite al usuario crear una cuenta nueva proporcionando
su nombre completo, correo electr\'onico, contrase\~na y confirmaci\'on de
contrase\~na. Incluye validaci\'on en tiempo real de todos los campos con el
esquema de validaci\'on Zod, que verifica que las contrase\~nas coincidan y que
los campos cumplan con los requisitos m\'inimos.

Despu\'es del registro exitoso, la aplicaci\'on muestra un mensaje de
confirmaci\'on y redirige autom\'aticamente al usuario hacia la pantalla de
inicio de sesi\'on.

\begin{figure}[H]
\centering
  \figuraAPA{fig:pantalla-registro}{Pantalla de registro de la aplicaci\'on \nombreApp{}}{%
    \includegraphics[width=0.35\linewidth]{pantalla-registro.png}%
  }{Elaboraci\'on propia (captura de pantalla de la aplicaci\'on).}
\end{figure}

Los campos del formulario de registro se detallan en la tabla.

\begin{table}[H]
  \centering
  \caption{Campos del formulario de registro}
  \contenidoConFuente{%
    \begin{tablaAPA}
      \begin{tabular}{@{}p{3.2cm}p{3.8cm}p{8.6cm}@{}}
        \toprule
        \textbf{Campo} & \textbf{Validaci\'on} & \textbf{Descripci\'on} \\
        \midrule
        Nombre completo & M\'inimo 3 caracteres & Campo de texto con
        capitalizaci\'on autom\'atica de palabras. \\
        \addlinespace
        Email & Formato v\'alido de correo & Campo de texto con tipo de
        teclado \textit{email-address}. \\
        \addlinespace
        Contrase\~na & M\'inimo 6 caracteres & Campo de entrada segura con
        autocompletado de nueva contrase\~na. \\
        \addlinespace
        Confirmar contrase\~na & Debe coincidir con contrase\~na & Campo de
        entrada segura que valida la igualdad con el campo anterior. \\
        \bottomrule
      \end{tabular}
    \end{tablaAPA}%
  }{Elaboraci\'on propia en base al c\'odigo fuente de la aplicaci\'on.}
\end{table}

\subsection{Pantalla principal (tabs)}

La pantalla principal de la aplicaci\'on se compone de una barra de
navegaci\'on inferior (\texttt{Tab Bar}) con cuatro pesta\~nas: Inicio, Buscar,
Favoritos y Perfil. Cada pesta\~na muestra un \'{i}cono de la librer\'ia
\texttt{@expo/vector-icons/Ionicons} y un t\'itulo descriptivo.

La pesta\~na de \textbf{Inicio} muestra el nombre de la aplicaci\'on y un
mensaje de bienvenida. La pesta\~na de \textbf{Buscar} presenta la
funcionalidad de b\'usqueda de escenarios. La pesta\~na de \textbf{Favoritos}
permite consultar los escenarios guardados. La pesta\~na de \textbf{Perfil}
muestra la informaci\'on del usuario autenticado con opci\'on de cerrar
sesi\'on.

\begin{figure}[H]
\centering
  \figuraAPA{fig:pantalla-tabs}{Pantalla principal con barra de navegaci\'on inferior}{%
    \includegraphics[width=0.55\linewidth]{pantalla-tabs.png}%
  }{Elaboraci\'on propia (captura de pantalla de la aplicaci\'on).}
\end{figure}

Los elementos de la barra de navegaci\'on se resumen en la tabla.

\begin{table}[H]
  \centering
  \caption{Pesta\~nas de la barra de navegaci\'on inferior}
  \contenidoConFuente{%
    \begin{tablaAPA}
      \begin{tabular}{@{}p{2.4cm}p{3.0cm}p{10.2cm}@{}}
        \toprule
        \textbf{Pesta\~na} & \textbf{\'Icono} & \textbf{Funcionalidad} \\
        \midrule
        Inicio & \textit{home-outline} & Pantalla principal con mensaje de
        bienvenida y exploraci\'on de escenarios. \\
        \addlinespace
        Buscar & \textit{search-outline} & B\'usqueda de escenarios deportivos
        por nombre o categor\'ia. \\
        \addlinespace
        Favoritos & \textit{star-outline} & Lista de escenarios guardados por
        el usuario. \\
        \addlinespace
        Perfil & \textit{person-outline} & Informaci\'on del usuario y
        opci\'on de cerrar sesi\'on. \\
        \bottomrule
      \end{tabular}
    \end{tablaAPA}%
  }{Elaboraci\'on propia en base al c\'odigo fuente de la aplicaci\'on.}
\end{table}

\subsection{Pantalla de detalle del escenario}

La pantalla de detalle del escenario recibe el identificador del escenario como
par\'ametro de ruta (\texttt{scenario/[id]}) y muestra la informaci\'on
asociada. En esta etapa se presenta una interfaz placeholder que muestra el
t\'itulo de la pantalla y el identificador del escenario, lo que permite validar
la navegaci\'on y la recepci\'on de par\'ametros entre pantallas.

\begin{figure}[H]
\centering
  \figuraAPA{fig:pantalla-scenario-detail}{Pantalla de detalle del escenario}{%
    \includegraphics[width=0.28\linewidth]{pantalla-scenario-detail.png}%
  }{Elaboraci\'on propia (captura de pantalla de la aplicaci\'on).}
\end{figure}

\subsection{Pantalla de perfil de usuario}

La pantalla de perfil muestra la informaci\'on del usuario autenticado,
incluyendo su nombre completo, correo electr\'onico y un bot\'on para cerrar
sesi\'on. El cierre de sesi\'on presenta un di\'alogo de confirmaci\'on antes de
ejecutar la acci\'on, y utiliza el contexto de autenticaci\'on
(\texttt{AuthContext}) para gestionar el estado de la sesi\'on.

\begin{figure}[H]
\centering
  \figuraAPA{fig:pantalla-perfil}{Pantalla de perfil de usuario}{%
    \includegraphics[width=0.35\linewidth]{pantalla-perfil.png}%
  }{Elaboraci\'on propia (captura de pantalla de la aplicaci\'on).}
\end{figure}
